\documentclass[11pt,a4paper]{article}

% ---------- Packages ----------
\usepackage[margin=1in]{geometry}
\usepackage{amsmath, amssymb, amsthm, mathtools}
\usepackage{microtype}
\usepackage[colorlinks=true,
            linkcolor=blue!50!black,
            urlcolor=blue!50!black,
            citecolor=blue!50!black]{hyperref}

% ---------- Theorem environments ----------
\theoremstyle{plain}
\newtheorem{theorem}{Theorem}[section]
\newtheorem{proposition}[theorem]{Proposition}
\newtheorem{lemma}[theorem]{Lemma}
\newtheorem{corollary}[theorem]{Corollary}

\theoremstyle{definition}
\newtheorem{definition}[theorem]{Definition}

\theoremstyle{remark}
\newtheorem{remark}[theorem]{Remark}

% ---------- Math shortcuts ----------
\newcommand{\R}{\mathbb{R}}
\newcommand{\E}{\mathbb{E}}
\newcommand{\F}{\mathcal{F}}
\newcommand{\Prob}{\mathbb{P}}
\newcommand{\Q}{\mathbb{Q}}
\newcommand{\dd}{\mathop{}\!\mathrm{d}}
\newcommand{\eqdef}{\overset{\mathrm{def}}{=}}
\newcommand{\BS}{\mathrm{BS}}
\newcommand{\Vega}{\mathcal{V}}

% ---------- Document metadata ----------
\title{The Greeks of European Options under Black--Scholes}
\author{Quant Finance Portfolio --- Phase 1, Algorithm 2}
\date{\today}

\begin{document}
\maketitle

\begin{abstract}
We derive the closed-form expressions for the five classical Greeks
\((\Delta, \Gamma, \Vega, \Theta, \rho)\) of European call and put
options under the Black--Scholes model. The derivations rest on a
single algebraic identity --- the equality of the two normal-density
factors that appear when differentiating the price formula --- which
is established once and reused throughout. We close with the
identities that link the Greeks: the Black--Scholes partial differential
equation, the Vega--Gamma identity, the conventions for Theta, and the
non-uniform convexity of the call price in volatility. The document
serves as a reference for the implementation in the project and for
the volatility-implied analysis that depends on Vega.
\end{abstract}

% ============================================================
\section{Setup}
% ============================================================

Throughout this document we work under the Black--Scholes model
introduced in the companion document \texttt{black\_scholes.tex}. We
fix \(S, K > 0\), \(r \in \R\), \(\sigma > 0\), and \(T > t \geq 0\),
and let \(\tau \eqdef T - t\). The price of a European call is
\begin{equation}
    C^{\BS}(S, K, r, \sigma, \tau)
    = S\, N(d_1) - K e^{-r\tau} N(d_2),
    \label{eq:bs}
\end{equation}
with
\begin{equation}
    d_1 = \frac{\log(S/K) + (r + \tfrac{1}{2}\sigma^2)\tau}{\sigma\sqrt{\tau}},
    \qquad
    d_2 = d_1 - \sigma\sqrt{\tau},
    \label{eq:d1d2}
\end{equation}
where \(N\) and \(\varphi\) denote the standard normal CDF and density
respectively, and the put price is
\begin{equation}
    P^{\BS} = K e^{-r\tau} N(-d_2) - S\, N(-d_1).
    \label{eq:put}
\end{equation}

A Greek is a partial derivative of the option price with respect to one
of its arguments. Each derivative is taken holding the remaining
arguments fixed.

\begin{definition}[The classical Greeks]
For a European option with price \(V \in \{C^{\BS}, P^{\BS}\}\), the
classical Greeks are
\[
    \Delta = \frac{\partial V}{\partial S},
    \qquad
    \Gamma = \frac{\partial^2 V}{\partial S^2},
    \qquad
    \Vega = \frac{\partial V}{\partial \sigma},
    \qquad
    \Theta = \frac{\partial V}{\partial t},
    \qquad
    \rho = \frac{\partial V}{\partial r}.
\]
\end{definition}

The Theta convention is the calendar-time derivative, not the
maturity-time derivative; the two differ by a sign because
\(\partial / \partial t = -\partial / \partial \tau\) when
\(\tau = T - t\) with \(T\) fixed. We use the calendar-time form
because it makes the Black--Scholes PDE
(Theorem~\ref{thm:bs-pde}) read without sign reversals.

\begin{remark}[Two readings of Greeks]
Greeks admit two complementary readings. As \emph{sensitivities},
they describe how the option price moves under marginal changes to
the inputs, and form the basis of the dynamic hedging used at every
options desk. As \emph{coefficients of a partial differential
equation}, they are the building blocks of the no-arbitrage condition
the price must satisfy: the Black--Scholes equation is a linear
relation among \(\Theta, \Gamma, \Delta\) and \(V\) itself. The
algebraic and the dynamical readings are not independent; they are
two manifestations of the same underlying replication argument.
\end{remark}

% ============================================================
\section{The Key Identity}
% ============================================================

The closed forms for the Greeks could in principle be obtained by
brute differentiation of~\eqref{eq:bs}. Direct computation produces a
profusion of terms involving \(\partial d_1 / \partial x\) and
\(\partial d_2 / \partial x\) for each input \(x\). Almost all of those
terms cancel, and the cancellation rests on a single identity which we
establish first.

\begin{lemma}[Density identity]
\label{lem:density-identity}
For all admissible \(S, K, r, \sigma, \tau\),
\begin{equation}
    S\, \varphi(d_1) = K e^{-r\tau}\, \varphi(d_2).
    \label{eq:density-identity}
\end{equation}
\end{lemma}

\begin{proof}
Compute the difference of the squares of \(d_1\) and \(d_2\), using
\(d_2 = d_1 - \sigma\sqrt{\tau}\):
\begin{align*}
    d_1^2 - d_2^2
    &= (d_1 - d_2)(d_1 + d_2)
    = \sigma\sqrt{\tau} \cdot \left(2 d_1 - \sigma\sqrt{\tau}\right) \\
    &= 2 d_1 \sigma\sqrt{\tau} - \sigma^2 \tau.
\end{align*}
Substitute the definition of \(d_1\) into the first term:
\[
    2 d_1 \sigma\sqrt{\tau}
    = 2 \log(S/K) + (2r + \sigma^2)\tau.
\]
Therefore
\[
    d_1^2 - d_2^2
    = 2\log(S/K) + 2 r \tau,
\]
which gives
\[
    \exp\!\left(\tfrac{1}{2}(d_2^2 - d_1^2)\right)
    = \frac{K e^{-r\tau}}{S}.
\]
The identity follows by multiplying by \(S\) and recognising the
density:
\[
    S\, \varphi(d_1)
    = S \cdot \frac{1}{\sqrt{2\pi}} e^{-d_1^2/2}
    = \frac{K e^{-r\tau}}{\sqrt{2\pi}} e^{-d_2^2/2}
    = K e^{-r\tau}\, \varphi(d_2). \qedhere
\]
\end{proof}

The identity holds because \(d_1\) and \(d_2\) are not independent; they
share the same numerator structure and differ only by \(\sigma\sqrt{\tau}\).
The symmetric quadratic form \(d_1^2 - d_2^2\) collapses to something
\(\sigma\)-free, so the imbalance between \(\varphi(d_1)\) and
\(\varphi(d_2)\) is exactly compensated by \(K e^{-r\tau}/S\).

\begin{remark}[Where the identity comes from]
The identity~\eqref{eq:density-identity} is not a coincidence.
It is the algebraic shadow of the change of num\'{e}raire used in the
derivation of the Black--Scholes formula. Recall from
\texttt{black\_scholes.tex} that \(N(d_1) = \Prob^S(S_T > K)\) under
the stock-num\'{e}raire measure, while
\(N(d_2) = \Q(S_T > K)\) under the bond-num\'{e}raire (risk-neutral)
measure. The two measures differ by the Radon--Nikodym derivative
\(d\Prob^S / d\Q = S_T B_t / (B_T S_t)\), and the densities of \(\log
S_T\) under the two measures shift by exactly the amount that equates
\(S\varphi(d_1)\) and \(K e^{-r\tau}\varphi(d_2)\) at the strike.
The cancellations we will exploit below are the algebraic image of
that probabilistic equivalence.
\end{remark}

% ============================================================
\section{The First-Order Greeks}
% ============================================================

We derive \(\Delta\), \(\Vega\), and \(\rho\) using
Lemma~\ref{lem:density-identity}. The strategy is identical in each
case: differentiate~\eqref{eq:bs}, isolate the
\(\varphi\)-terms, and apply~\eqref{eq:density-identity} to make them
cancel. The remaining terms give the closed form.

% ------------------------------------------------------------
\subsection{Delta}
% ------------------------------------------------------------

\begin{proposition}[Delta]
\label{prop:delta}
\[
    \Delta_C = N(d_1), \qquad \Delta_P = N(d_1) - 1 = -N(-d_1).
\]
\end{proposition}

\begin{proof}
Differentiate~\eqref{eq:bs} with respect to \(S\):
\[
    \Delta_C = N(d_1) + S \varphi(d_1) \frac{\partial d_1}{\partial S}
              - K e^{-r\tau} \varphi(d_2) \frac{\partial d_2}{\partial S}.
\]
From~\eqref{eq:d1d2}, both \(\partial d_1 / \partial S\) and
\(\partial d_2 / \partial S\) equal \(1 / (S \sigma \sqrt{\tau})\)
(since \(d_2 - d_1\) does not depend on \(S\)). Thus
\[
    \Delta_C = N(d_1) + \frac{1}{S \sigma \sqrt{\tau}}
        \bigl[S \varphi(d_1) - K e^{-r\tau} \varphi(d_2)\bigr],
\]
and the bracketed term vanishes by~\eqref{eq:density-identity}.
Therefore \(\Delta_C = N(d_1)\). For the put, differentiate
\eqref{eq:put} or use put-call parity \(C - P = S - Ke^{-r\tau}\),
which gives \(\Delta_C - \Delta_P = 1\).
\end{proof}

\begin{remark}[Replication interpretation]
The Black--Scholes price~\eqref{eq:bs} is structured as
\(S \cdot N(d_1) - K e^{-r\tau} \cdot N(d_2)\), which is the value of
a portfolio holding \(N(d_1)\) shares financed by borrowing
\(K e^{-r\tau} N(d_2)\) at the riskless rate. Proposition~\ref{prop:delta}
makes this visible: the share quantity \(N(d_1)\) is exactly the
hedge ratio, the option's sensitivity to the underlying. The fact
that \(\partial C/\partial S\) lands cleanly on \(N(d_1)\) ---
without trailing \(\varphi\) terms --- is not algebraic luck, but the
internal consistency of the replication argument that produced the
formula in the first place.
\end{remark}

% ------------------------------------------------------------
\subsection{Vega}
% ------------------------------------------------------------

\begin{proposition}[Vega]
\label{prop:vega}
\[
    \Vega_C = \Vega_P = S \varphi(d_1) \sqrt{\tau}.
\]
\end{proposition}

\begin{proof}
Differentiate~\eqref{eq:bs} with respect to \(\sigma\):
\[
    \Vega_C = S \varphi(d_1) \frac{\partial d_1}{\partial \sigma}
            - K e^{-r\tau} \varphi(d_2) \frac{\partial d_2}{\partial \sigma}.
\]
By~\eqref{eq:density-identity}, the two density factors satisfy
\(S\varphi(d_1) = K e^{-r\tau} \varphi(d_2)\), so this collapses to
\[
    \Vega_C = S \varphi(d_1)\!\left(
        \frac{\partial d_1}{\partial \sigma}
        - \frac{\partial d_2}{\partial \sigma}
    \right).
\]
From \(d_1 - d_2 = \sigma\sqrt{\tau}\),
\(\partial d_1/\partial\sigma - \partial d_2/\partial\sigma = \sqrt{\tau}\),
giving \(\Vega_C = S \varphi(d_1) \sqrt{\tau}\).

For the put, differentiate~\eqref{eq:put} or apply put-call parity:
\(\partial(C - P)/\partial\sigma = 0\), since the parity right-hand
side does not depend on \(\sigma\). Hence \(\Vega_P = \Vega_C\).
\end{proof}

\begin{remark}[Vega is positive, and identical for calls and puts]
Vega is strictly positive because all three factors are. This is the
monotonicity that underlies the well-posedness of the
implied-volatility inversion (see \texttt{implied\_volatility.tex},
Lemma 2.1). The equality \(\Vega_C = \Vega_P\) is not a coincidence
either: parity differentiated in \(\sigma\) has zero on the right.
The same argument gives \(\Gamma_C = \Gamma_P\) and
\(\Theta_C - \Theta_P = -r K e^{-r\tau}\); only Delta and Rho carry
genuine option-type information.
\end{remark}

% ------------------------------------------------------------
\subsection{Rho}
% ------------------------------------------------------------

\begin{proposition}[Rho]
\label{prop:rho}
\[
    \rho_C = K \tau e^{-r\tau} N(d_2),
    \qquad
    \rho_P = -K \tau e^{-r\tau} N(-d_2).
\]
\end{proposition}

\begin{proof}
Differentiate~\eqref{eq:bs} with respect to \(r\):
\[
    \rho_C
    = S \varphi(d_1) \frac{\partial d_1}{\partial r}
        + K \tau e^{-r\tau} N(d_2)
        - K e^{-r\tau} \varphi(d_2) \frac{\partial d_2}{\partial r}.
\]
Both \(\partial d_1 / \partial r\) and \(\partial d_2 / \partial r\)
equal \(\sqrt{\tau}/\sigma\) (a constant independent of \(d_1\) versus
\(d_2\) shift), so
\(\bigl[S\varphi(d_1) - K e^{-r\tau}\varphi(d_2)\bigr] \cdot \sqrt{\tau}/\sigma\)
vanishes again by~\eqref{eq:density-identity}. Only the
\(K\tau e^{-r\tau} N(d_2)\) term survives. The put result follows from
parity.
\end{proof}

% ============================================================
\section{The Second-Order and Time Greeks}
% ============================================================

The remaining Greeks --- \(\Gamma\) and \(\Theta\) --- arise from
second derivatives or from time. \(\Gamma\) gets a separate
treatment because it does not require~\eqref{eq:density-identity}; it
falls out of differentiating \(\Delta = N(d_1)\) directly. \(\Theta\)
involves the most cancellation and is the longest derivation.

% ------------------------------------------------------------
\subsection{Gamma}
% ------------------------------------------------------------

\begin{proposition}[Gamma]
\label{prop:gamma}
\[
    \Gamma_C = \Gamma_P = \frac{\varphi(d_1)}{S \sigma \sqrt{\tau}}.
\]
\end{proposition}

\begin{proof}
By Proposition~\ref{prop:delta}, \(\Delta_C = N(d_1)\). Differentiating
in \(S\):
\[
    \Gamma_C = \varphi(d_1) \frac{\partial d_1}{\partial S}
             = \frac{\varphi(d_1)}{S \sigma \sqrt{\tau}}.
\]
The put case follows from \(\Delta_C - \Delta_P = 1\) (constant in
\(S\)).
\end{proof}

% ------------------------------------------------------------
\subsection{Theta}
% ------------------------------------------------------------

\begin{proposition}[Theta]
\label{prop:theta}
\begin{align*}
    \Theta_C
    &= -\frac{S \varphi(d_1) \sigma}{2 \sqrt{\tau}}
        - r K e^{-r\tau} N(d_2), \\
    \Theta_P
    &= -\frac{S \varphi(d_1) \sigma}{2 \sqrt{\tau}}
        + r K e^{-r\tau} N(-d_2).
\end{align*}
\end{proposition}

\begin{proof}
Recall \(\tau = T - t\), so \(\partial/\partial t = -\partial/\partial\tau\).
Differentiate~\eqref{eq:bs} with respect to \(\tau\):
\[
    \frac{\partial C^{\BS}}{\partial \tau}
    = S \varphi(d_1) \frac{\partial d_1}{\partial \tau}
        + K r e^{-r\tau} N(d_2)
        - K e^{-r\tau} \varphi(d_2) \frac{\partial d_2}{\partial \tau}.
\]
The second term is straightforward. For the first and third,
\[
    \frac{\partial d_1}{\partial \tau}
    = \frac{r + \tfrac{1}{2}\sigma^2}{\sigma\sqrt{\tau}}
        - \frac{d_1}{2\tau},
    \qquad
    \frac{\partial d_2}{\partial \tau}
    = \frac{r - \tfrac{1}{2}\sigma^2}{\sigma\sqrt{\tau}}
        - \frac{d_2}{2\tau}.
\]
Group the \(\varphi\)-terms. By Lemma~\ref{lem:density-identity},
\(S\varphi(d_1) = K e^{-r\tau}\varphi(d_2)\), so we factor out a single
\(S\varphi(d_1)\) from both:
\[
    S\varphi(d_1)\!\left[
        \frac{\partial d_1}{\partial\tau} - \frac{\partial d_2}{\partial\tau}
    \right]
    = S\varphi(d_1)\!\left[
        \frac{\sigma}{\sqrt{\tau}}
        + \frac{d_2 - d_1}{2\tau}
    \right].
\]
Since \(d_1 - d_2 = \sigma\sqrt{\tau}\), the bracketed expression
becomes
\[
    \frac{\sigma}{\sqrt{\tau}} - \frac{\sigma}{2\sqrt{\tau}}
    = \frac{\sigma}{2\sqrt{\tau}}.
\]
Therefore
\[
    \frac{\partial C^{\BS}}{\partial \tau}
    = \frac{S\varphi(d_1) \sigma}{2 \sqrt{\tau}} + r K e^{-r\tau} N(d_2),
\]
and the calendar-time derivative is the negative,
\[
    \Theta_C
    = \frac{\partial C^{\BS}}{\partial t}
    = -\frac{\partial C^{\BS}}{\partial \tau}
    = -\frac{S\varphi(d_1) \sigma}{2 \sqrt{\tau}} - r K e^{-r\tau} N(d_2).
\]
The put expression follows from parity differentiated in \(t\):
\(\Theta_C - \Theta_P = \partial(C - P)/\partial t
= \partial(S - K e^{-r\tau})/\partial t = -r K e^{-r\tau}\).
\end{proof}

\begin{remark}[Sign of Theta]
For a non-dividend-paying call, the leading \(\varphi\)-term is always
negative, and the discount term is also negative whenever
\(N(d_2) > 0\). Hence \(\Theta_C < 0\): the call loses value as time
passes, all else equal. For the put, the sign of \(\Theta_P\) is
ambiguous because the two terms have opposite signs; deep ITM puts
near expiry can have positive Theta when the discount benefit
outweighs the time-value decay.
\end{remark}

% ============================================================
\section{Identities and Structural Properties}
% ============================================================

The closed forms hide several relations among the Greeks. We collect
them here because they appear repeatedly in pricing, hedging, and
validation.

% ------------------------------------------------------------
\subsection{The Black--Scholes PDE in Greeks}
% ------------------------------------------------------------

\begin{theorem}[Black--Scholes PDE]
\label{thm:bs-pde}
The price of a European call (or put) under the Black--Scholes model
satisfies
\begin{equation}
    \Theta + \tfrac{1}{2}\sigma^2 S^2\, \Gamma + r S\, \Delta - r V = 0,
    \label{eq:bs-pde}
\end{equation}
identically in \((S, K, r, \sigma, \tau)\), with \(V \in \{C^{\BS}, P^{\BS}\}\).
\end{theorem}

\begin{proof}
Substituting the closed forms from
Propositions~\ref{prop:delta},~\ref{prop:gamma}, and~\ref{prop:theta}
into~\eqref{eq:bs-pde} for the call:
\begin{align*}
&-\frac{S\varphi(d_1)\sigma}{2\sqrt{\tau}} - rKe^{-r\tau} N(d_2) \\
&\qquad + \tfrac{1}{2}\sigma^2 S^2 \cdot \frac{\varphi(d_1)}{S\sigma\sqrt{\tau}}
    + r S\, N(d_1) - r\bigl[S N(d_1) - K e^{-r\tau} N(d_2)\bigr] \\
&= -\frac{S\varphi(d_1)\sigma}{2\sqrt{\tau}} + \frac{S\varphi(d_1)\sigma}{2\sqrt{\tau}}
    + r S N(d_1) - r S N(d_1) - r K e^{-r\tau} N(d_2) + r K e^{-r\tau} N(d_2) \\
&= 0.
\end{align*}
The put case follows by parity, or by the same substitution with the
put closed forms.
\end{proof}

\begin{remark}[Why this identity matters]
Theorem~\ref{thm:bs-pde} couples four of the five classical Greeks
(\(\Vega\) and \(\rho\) do not appear because the model treats
\(\sigma\) and \(r\) as constants). The PDE is therefore a stringent
joint test of any closed-form or numerical implementation: a sign
error or missing factor in a single Greek breaks the residual at
machine precision. The numerical validation suite in
\texttt{validate\_greeks.py} exploits this directly to triangulate
correctness.
\end{remark}

% ------------------------------------------------------------
\subsection{The Vega--Gamma identity}
% ------------------------------------------------------------

\begin{proposition}[Vega--Gamma identity]
\label{prop:vega-gamma}
\begin{equation}
    \Vega = S^2\, \sigma\, \tau\, \Gamma.
    \label{eq:vega-gamma}
\end{equation}
\end{proposition}

\begin{proof}
Substituting the closed forms,
\(S^2 \sigma \tau \cdot \varphi(d_1) / (S \sigma \sqrt{\tau}) = S \varphi(d_1) \sqrt{\tau} = \Vega\).
\end{proof}

\begin{remark}[The dual nature of Vega and Gamma]
The Vega--Gamma identity says that, for a vanilla European option,
being long Vega and being long Gamma are the same exposure up to a
deterministic rescaling. This collapses two notions of risk that the
trader's intuition might keep separate. The collapse is specific to
Black--Scholes: in stochastic volatility models (Heston, SABR),
\(\sigma\) becomes a state variable in its own right, the
Vega-equivalent splits into multiple sensitivities to the volatility
process, and~\eqref{eq:vega-gamma} no longer holds. The Black--Scholes
collapse is one more way in which the model is a degenerate special
case of the broader theory.
\end{remark}

% ------------------------------------------------------------
\subsection{Curvature in volatility}
% ------------------------------------------------------------

The first derivative of \(C^{\BS}\) in \(\sigma\) (Vega) is strictly
positive. The second derivative (sometimes called \emph{Volga}) is
not.

\begin{proposition}[Volga]
\label{prop:volga}
\begin{equation}
    \frac{\partial^2 C^{\BS}}{\partial \sigma^2}
    = \frac{S \varphi(d_1) \sqrt{\tau}}{\sigma}\, d_1\, d_2.
    \label{eq:volga}
\end{equation}
\end{proposition}

\begin{proof}
Differentiate \(\Vega = S\varphi(d_1)\sqrt{\tau}\) with respect to
\(\sigma\):
\[
    \frac{\partial \Vega}{\partial \sigma}
    = S \sqrt{\tau} \cdot \varphi'(d_1) \cdot \frac{\partial d_1}{\partial \sigma}.
\]
Using \(\varphi'(x) = -x \varphi(x)\) and
\(\partial d_1 / \partial \sigma = -d_2 / \sigma\) (a calculation that
isolates the explicit \(\sigma\)-dependence and uses
\(d_2 = d_1 - \sigma\sqrt{\tau}\)):
\[
    \frac{\partial \Vega}{\partial \sigma}
    = S \sqrt{\tau} \cdot (-d_1)\varphi(d_1) \cdot \frac{-d_2}{\sigma}
    = \frac{S\varphi(d_1)\sqrt{\tau}}{\sigma}\, d_1 d_2. \qedhere
\]
\end{proof}

The sign of Volga therefore tracks the sign of \(d_1 d_2\). For
options moderately near at-the-money with small-to-moderate \(\sigma\),
\(d_1\) and \(d_2\) typically straddle zero and the call is locally
\emph{concave} in \(\sigma\); for options far from ATM or for large
\(\sigma\), they share signs and the call is convex. The change of
curvature occurs along the locus \(d_1 d_2 = 0\), which is the central
object of the smile-implied-volatility approximation literature.

\begin{remark}[Implication for Newton's method on implied volatility]
Volga's sign change is the technical reason why monotonicity of
\(C^{\BS}\) in \(\sigma\) is not by itself sufficient for global
convergence of Newton's method when computing implied volatility, as
discussed in \texttt{implied\_volatility.tex} (Section 3.2). When the
initial guess and the true root lie on opposite sides of the locus
\(d_1 d_2 = 0\), the first Newton step can overshoot.
This is why the implementation uses a region-aware initial guess
(Brenner--Subrahmanyam near ATM, neutral default elsewhere) rather
than a constant.
\end{remark}

% ============================================================
\section{Summary}
% ============================================================

\begin{center}
\renewcommand{\arraystretch}{1.5}
\begin{tabular}{l|l|l}
Greek & Call & Put \\\hline
\(\Delta\) & \(N(d_1)\) & \(N(d_1) - 1\) \\
\(\Gamma\) & \multicolumn{2}{l}{\(\dfrac{\varphi(d_1)}{S\sigma\sqrt{\tau}}\)} \\[0.5em]
\(\Vega\) & \multicolumn{2}{l}{\(S \varphi(d_1) \sqrt{\tau}\)} \\[0.5em]
\(\Theta\)
    & \(-\dfrac{S\varphi(d_1)\sigma}{2\sqrt{\tau}} - rKe^{-r\tau}N(d_2)\)
    & \(-\dfrac{S\varphi(d_1)\sigma}{2\sqrt{\tau}} + rKe^{-r\tau}N(-d_2)\) \\[0.7em]
\(\rho\) & \(K\tau e^{-r\tau} N(d_2)\) & \(-K\tau e^{-r\tau} N(-d_2)\)
\end{tabular}
\end{center}

The Greeks satisfy three structural identities: the Black--Scholes
PDE \(\Theta + \tfrac{1}{2}\sigma^2 S^2 \Gamma + rS\Delta - rV = 0\),
the Vega--Gamma identity \(\Vega = S^2 \sigma \tau \Gamma\), and the
non-trivial sign of Volga \(\partial \Vega / \partial \sigma\) which
governs the convexity of the price in volatility. Each identity is
exploited downstream: the PDE in the validation of the Greeks
implementation, the Vega--Gamma identity as an algebraic sanity
check, and Volga in the analysis of Newton's method for implied
volatility.

\end{document}
